% Part: first-order-logic
% Chapter: natural-deduction
% Section: proving-things

\documentclass[../../../include/open-logic-section]{subfiles}

\begin{document}

\iftag{FOL}
      {\olfileid{fol}{ntd}{pro}}
      {\olfileid{pl}{ntd}{pro}}

\olsection{Tuladha \usetoken{P}{derivation}}

\begin{ex}
Ayo menehi !!a{derivation} saka !!{sentence} $(!A \land !B) \lif !A$.

Kita miwiti kanthi nulis kesimpulan kang dikarepake ing sisih ngisor
!!{derivation}.
\begin{prooftree}
\AxiomC{}
\UnaryInfC{$(!A\land !B) \lif !A$}
\end{prooftree}

Sabanjure, kita kudu nemtokake jinis inferensi apa kang bisa ngasilake
!!a{sentence} kanthi wujud iki. !!{main operator} saka kesimpulan iku
$\lif$, mula kita bakal nyoba nggayuh kesimpulan nganggo aturan
\Intro{\lif}. Becik yen asumsi kang digunakake lan label aturan
inferensi ditulis sairing lumakune pambuktèn, supaya ing pungkasan
bukti gampang dipriksa apa kabeh asumsi wis !!{discharged}.
\begin{prooftree}
\AxiomC{$\Discharge{!A \land !B}{1}$}
\DeduceC{$!A$}
\DischargeRule{\Intro{\lif}}{1}
\UnaryInfC{$(!A\land !B) \lif !A$}
\end{prooftree}

Saiki kita kudu ngisi langkah-langkah saka asumsi $!A \land !B$
menyang $!A$. Amarga mung ana siji konektif kang kudu ditangani, yaiku
$\land$, kita kudu nggunakake aturan eliminasi $\land$. Kanthi mangkono
kita entuk bukti iki:
\begin{prooftree}
\AxiomC{$\Discharge{!A \land !B}{1}$}
\RightLabel{\Elim{\land}}
\UnaryInfC{$!A$}
\DischargeRule{\Intro{\lif}}{1}
\UnaryInfC{$(!A\land !B) \lif !A$}
\end{prooftree}
Saiki kita duwe !!{derivation} kang bener saka $(!A \land
!B) \lif !A$.
\end{ex}

\begin{ex}
Saiki ayo menehi !!a{derivation} saka $(\lnot !A \lor !B)
\lif (!A \lif !B)$.

Kita miwiti kanthi nulis kesimpulan kang dikarepake ing sisih ngisor
!!{derivation}.
\begin{prooftree}
\AxiomC{}
\UnaryInfC{$(\lnot !A \lor !B) \lif (!A \lif !B)$}
\end{prooftree}
Kanggo nemokake aturan logis kang bisa menehi kesimpulan iki, kita
ndeleng konektif logis ing kesimpulan: $\lnot$, $\lor$, lan $\lif$.
Saiki kita mung nggatekake kedadeyan kapisan saka $\lif$, amarga iku
!!{main operator} saka !!{sentence} pungkasan, dene $\lnot$, $\lor$,
lan kedadeyan kapindho saka $\lif$ ana ing jangkauan konektif liyane,
mula bakal ditangani mengko. Dadi, kita miwiti nganggo aturan
\Intro{\lif}. Panrapan kang bener kudu awujud mangkene:
\begin{prooftree}
\AxiomC{$\Discharge{\lnot !A \lor !B}{1}$}
\DeduceC{$!A \lif !B$}
\DischargeRule{\Intro{\lif}}{1}
\UnaryInfC{$(\lnot !A \lor !B) \lif (!A \lif !B)$}
\end{prooftree}
Iki menehi rong pilihan kanggo nerusake. Kita bisa terus makarya saka
ngisor menyang ndhuwur lan nggolèki panrapan aturan \Intro{\lif}
liyane, utawa makarya saka ndhuwur menyang ngisor lan ngetrapake
aturan \Elim{\lor}. Ayo ngetrapake pilihan kapindho. Kita bakal
nggunakake asumsi $\lnot !A \lor !B$ minangka premis paling kiwa saka
\Elim{\lor}. Kanggo panrapan \Elim{\lor} kang bener, premis loro
liyane kudu padha karo kesimpulan $!A \lif !B$, nanging saben premis
bisa diasilake kanthi urut saka asumsi liyane, yaiku salah siji saka
rong disjung ing $\lnot !A \lor !B$. Mula !!{derivation} kita bakal
awujud mangkene:
\begin{prooftree}
\AxiomC{$\Discharge{\lnot !A \lor !B}{1}$}
\AxiomC{$\Discharge{\lnot !A}{2}$}
\DeduceC{$!A \lif !B$}
\AxiomC{$\Discharge{!B}{2}$}
\DeduceC{$!A \lif !B$}
\DischargeRule{\Elim{\lor}}{2}
\TrinaryInfC{$!A \lif !B$}
\DischargeRule{\Intro{\lif}}{1}
\UnaryInfC{$(\lnot !A \lor !B) \lif (!A \lif !B)$}
\end{prooftree}

Ing saben cabang loro ing sisih tengen, kita arep !!{derive} $!A
\lif !B$, kang paling becik ditindakake nganggo \Intro{\lif}.
\begin{prooftree}
\AxiomC{$\Discharge{\lnot !A \lor !B}{1}$}
\AxiomC{$\Discharge{\lnot !A}{2}, \Discharge{!A}{3}$}
\DeduceC{$!B$}
\DischargeRule{\Intro{\lif}}{3}
\UnaryInfC{$!A \lif !B$}
\AxiomC{$\Discharge{!B}{2}, \Discharge{!A}{4}$}
\DeduceC{$!B$}
\DischargeRule{\Intro{\lif}}{4}
\UnaryInfC{$!A \lif !B$}
\DischargeRule{\Elim{\lor}}{2}
\TrinaryInfC{$!A \lif !B$}
\DischargeRule{\Intro{\lif}}{1}
\UnaryInfC{$(\lnot !A \lor !B) \lif (!A \lif !B)$}
\end{prooftree}

Kanggo rong bagean !!{derivation} kang isih kosong, kita mbutuhake
!!{derivation}s saka $!B$ adhedhasar $\lnot !A$ lan $!A$ ing cabang
tengah, lan adhedhasar $!A$ lan $!B$ ing cabang tengen. Cathetan
koreksi sumber OLPL-011: teks Inggris nyebut cabang kapindho iki
minangka cabang kiwa, nanging wit lan paragraf sabanjure kanthi cetha
ndunungake ing sisih tengen. Ayo luwih dhisik njupuk cabang tengah.
$\lnot !A$ lan $!A$ iku rong premis saka \Elim{\lnot}:
\begin{prooftree}
\AxiomC{$\Discharge{\lnot !A}{2}$}
\AxiomC{$\Discharge{!A}{3}$}
\RightLabel{\Elim{\lnot}}
\BinaryInfC{$\lfalse$}
\DeduceC{$!B$}
\end{prooftree}
Kanthi nggunakake \FalseInt, kita bisa ngasilake $!B$ minangka
kesimpulan lan ngrampungake cabang iki.
\begin{prooftree}
\AxiomC{$\Discharge{\lnot !A \lor !B}{1}$}
\AxiomC{$\Discharge{\lnot !A}{2}$}
\AxiomC{$\Discharge{!A}{3}$}
\RightLabel{\Elim{\lnot}}
\BinaryInfC{$\lfalse$}
\RightLabel{\FalseInt}
\UnaryInfC{$!B$}
\DischargeRule{\Intro{\lif}}{3}
\UnaryInfC{$!A \lif !B$}
\AxiomC{$\Discharge{!B}{2}, \Discharge{!A}{4}$}
\DeduceC{$!B$}
\DischargeRule{\Intro{\lif}}{4}
\UnaryInfC{$!A \lif !B$}
\DischargeRule{\Elim{\lor}}{2}
\TrinaryInfC{$!A \lif !B$}
\DischargeRule{\Intro{\lif}}{1}
\UnaryInfC{$(\lnot !A \lor !B) \lif (!A \lif !B)$}
\end{prooftree}
Cathetan koreksi sumber OLPL-010: label ing wit Inggris ngarani
langkah saka negasi lan ukara menyang kontradiksi minangka introduksi
falsum, nanging wit parsial sadurunge lan aturan kang ditetepake
ngenali langkah iki minangka eliminasi negasi.

Saiki ayo ndeleng cabang paling tengen. Ing kene kudu dielingi yen
definisi !!{derivation} \emph{ngidini asumsi dibubarake}, nanging
\emph{ora meksa} asumsi kasebut dibubarake. Kanthi tembung liya, yen
kita bisa ngasilake $!B$ saka salah siji asumsi $!A$ lan $!B$ tanpa
nggunakake liyane, iku ora dadi masalah. Lan kanggo !!{derive} $!B$
saka~$!B$ iku gampang: $!B$ dhewe wis dadi !!a{derivation} kaya
mangkono, tanpa mbutuhake inferensi. Dadi, kita cukup mbusak
asumsi~$!A$.
\begin{prooftree}
\AxiomC{$\Discharge{\lnot !A \lor !B}{1}$}
\AxiomC{$\Discharge{\lnot !A}{2}$}
\AxiomC{$\Discharge{!A}{3}$}
\RightLabel{\Elim{\lnot}}
\BinaryInfC{$\lfalse$}
\RightLabel{\FalseInt}
\UnaryInfC{$!B$}
\DischargeRule{\Intro{\lif}}{3}
\UnaryInfC{$!A \lif !B$}
\AxiomC{$\Discharge{!B}{2}$}
\RightLabel{\Intro{\lif}}
\UnaryInfC{$!A \lif !B$}
\DischargeRule{\Elim{\lor}}{2}
\TrinaryInfC{$!A \lif !B$}
\DischargeRule{\Intro{\lif}}{1}
\UnaryInfC{$(\lnot !A \lor !B) \lif (!A \lif !B)$}
\end{prooftree}
Gatekna yen ing !!{derivation} kang wis rampung, inferensi
\Intro{\lif} paling tengen sejatine ora mbubarake asumsi apa wae.
\end{ex}

\begin{ex}
Nganti saiki kita durung mbutuhake aturan \FalseCl{}. Aturan iki
mligi amarga ngidini kita mbubarake asumsi kang dudu
sub-!!{formula} saka kesimpulan aturan kasebut. Aturan iki raket
sesambungane karo aturan \FalseInt{}. Malah, aturan \FalseInt{} iku
kasus mligi saka aturan \FalseCl{}---ana logika kang diarani
``logika intuisionistik'' kang mung ngidini \FalseInt{}. Aturan
\FalseCl{} dadi upaya pungkasan nalika ora ana cara liyane kang bisa
ditindakake. Tuladhane, upamana kita arep !!{derive} $!A \lor \lnot
!A$. Strategi lumrahe yaiku nyoba !!{derive} $!A \lor \lnot !A$
nganggo $\Intro{\lor}$. Nanging iki mbutuhake supaya kita !!{derive}
$!A$ utawa $\lnot !A$ tanpa asumsi, lan prakara iki ora bisa
ditindakake. \FalseCl{} bisa nulungi!
\begin{prooftree}
  \AxiomC{$\Discharge{\lnot(!A \lor \lnot !A)}{1}$}
  \DeduceC{$\lfalse$}
  \DischargeRule{\FalseCl}{1}
  \UnaryInfC{$!A \lor \lnot !A$}
\end{prooftree}
Saiki kita nggolèki !!a{derivation} saka $\lfalse$ adhedhasar
$\lnot(!A \lor \lnot !A)$. Amarga $\lfalse$ iku kesimpulan saka
$\Elim{\lnot}$, kita bisa nyoba aturan kasebut:
\begin{prooftree}
  \AxiomC{$\Discharge{\lnot(!A \lor \lnot !A)}{1}$}
  \DeduceC{$\lnot !A$}
  \AxiomC{$\Discharge{\lnot(!A \lor \lnot !A)}{1}$}
  \DeduceC{$!A$}
  \RightLabel{\Elim{\lnot}}
  \BinaryInfC{$\lfalse$}
  \DischargeRule{\FalseCl}{1}
  \UnaryInfC{$!A \lor \lnot !A$}
\end{prooftree}
Strategi kanggo nemokake !!a{derivation} saka~$\lnot !A$ mbutuhake
panrapan~$\Intro{\lnot}$:
\begin{prooftree}
  \AxiomC{$\Discharge{\lnot(!A \lor \lnot !A)}{1}, \Discharge{!A}{2}$}
  \DeduceC{$\lfalse$}
  \DischargeRule{\Intro{\lnot}}{2}
  \UnaryInfC{$\lnot !A$}
  \AxiomC{$\Discharge{\lnot(!A \lor \lnot !A)}{1}$}
  \DeduceC{$!A$}
  \RightLabel{\Elim{\lnot}}
  \BinaryInfC{$\lfalse$}
  \DischargeRule{\FalseCl}{1}
  \UnaryInfC{$!A \lor \lnot !A$}
\end{prooftree}
Ing kene, kita bisa ngasilake $\lfalse$ kanthi gampang kanthi
ngetrapake $\Elim{\lnot}$ marang asumsi $\lnot(!A \lor \lnot !A)$ lan
$!A \lor \lnot !A$, kang ndherek saka asumsi anyar $!A$
liwat~$\Intro{\lor}$:
\begin{prooftree}
  \AxiomC{$\Discharge{\lnot(!A \lor \lnot !A)}{1}$}
  \AxiomC{$\Discharge{!A}{2}$}
  \RightLabel{\Intro{\lor}}
  \UnaryInfC{$!A \lor \lnot !A$}
  \RightLabel{\Elim{\lnot}}
  \BinaryInfC{$\lfalse$}
  \DischargeRule{\Intro{\lnot}}{2}
  \UnaryInfC{$\lnot !A$}
  \AxiomC{$\Discharge{\lnot(!A \lor \lnot !A)}{1}$}
  \DeduceC{$!A$}
  \RightLabel{\Elim{\lnot}}
  \BinaryInfC{$\lfalse$}
  \DischargeRule{\FalseCl}{1}
  \UnaryInfC{$!A \lor \lnot !A$}
\end{prooftree}
Ing sisih tengen kita nggunakake strategi kang padha, kajaba $!A$
diasilake liwat~\FalseCl:
\begin{prooftree}
  \AxiomC{$\Discharge{\lnot(!A \lor \lnot !A)}{1}$}
  \AxiomC{$\Discharge{!A}{2}$}
  \RightLabel{\Intro{\lor}}
  \UnaryInfC{$!A \lor \lnot !A$}
  \RightLabel{\Elim{\lnot}}
  \BinaryInfC{$\lfalse$}
  \DischargeRule{\Intro{\lnot}}{2}
  \UnaryInfC{$\lnot !A$}
  \AxiomC{$\Discharge{\lnot(!A \lor \lnot !A)}{1}$}
  \AxiomC{$\Discharge{\lnot !A}{3}$}
  \RightLabel{\Intro{\lor}}
  \UnaryInfC{$!A \lor \lnot !A$}
  \RightLabel{\Elim{\lnot}}
  \BinaryInfC{$\lfalse$}
  \DischargeRule{\FalseCl}{3}
  \UnaryInfC{$!A$}
  \RightLabel{\Elim{\lnot}}
  \BinaryInfC{$\lfalse$}
  \DischargeRule{\FalseCl}{1}
  \UnaryInfC{$!A \lor \lnot !A$}
\end{prooftree}
\end{ex}

\begin{prob}
Wènèhana !!{derivation}s kang nuduhake prakara-prakara iki:
\begin{enumerate}
\item $!A \land (!B \land !C) \Proves (!A \land !B) \land !C$.
\item $!A \lor (!B \lor !C) \Proves (!A \lor !B) \lor !C$.
\item $!A \lif (!B \lif !C) \Proves !B \lif (!A \lif !C)$.
\item $!A \Proves \lnot\lnot !A$.
\end{enumerate}
\end{prob}

\begin{prob}
Wènèhana !!{derivation}s kang nuduhake prakara-prakara iki:
\begin{enumerate}
\item $(!A \lor !B) \lif !C \Proves !A \lif !C$.
\item $(!A \lif !C) \land (!B \lif !C) \Proves (!A \lor !B) \lif !C$.
\item $\Proves \lnot(!A \land \lnot !A)$.
\item $!B \lif !A \Proves \lnot !A \lif \lnot !B$.
\item $\Proves (!A \lif \lnot !A) \lif \lnot !A$.
\item $\Proves \lnot(!A \lif !B) \lif \lnot !B$.
\item $!A \lif !C \Proves \lnot (!A \land \lnot !C)$.
\item $!A \land \lnot !C \Proves \lnot (!A \lif !C)$.
\item $!A \lor !B, \lnot !B \Proves !A$.
\item $\lnot !A \lor \lnot !B \Proves \lnot(!A \land !B)$.
\item $\Proves (\lnot !A \land \lnot !B) \lif\lnot(!A \lor !B)$.
\item $\Proves \lnot(!A \lor !B) \lif (\lnot !A \land \lnot !B)$.
\end{enumerate}
\end{prob}

\begin{prob}
Wènèhana !!{derivation}s kang nuduhake prakara-prakara iki:
\begin{enumerate}
\item $\lnot(!A \lif !B) \Proves !A$.
\item $\lnot(!A \land !B) \Proves \lnot !A \lor \lnot !B$.
\item $!A \lif !B \Proves \lnot !A \lor !B$.
\item $\Proves \lnot \lnot !A \lif !A$.
\item $!A \lif !B, \lnot !A \lif !B \Proves !B$.
\item $(!A \land !B) \lif !C \Proves (!A \lif !C) \lor (!B \lif !C)$.
\item $(!A \lif !B) \lif !A \Proves !A$.
\item $\Proves (!A \lif !B) \lor (!B \lif !C)$.
\end{enumerate}
(Kabeh iki mbutuhake aturan~$\FalseCl$.)
\end{prob}

\end{document}
